\section{Related Work}
\label{sec:related}

\textbf{Exploit-likelihood prioritization.} EPSS~\cite{jacobs2021epss} reframed prioritization
from static severity (CVSS)~\cite{cvss31} toward predicted exploitation, replacing a question
about how bad a vulnerability could be in principle with a question about how likely it is to be
exploited in fact. Subsequent data-driven work refined exploit prediction with richer signals and
community-sourced features~\cite{jacobs2023enhancing}, and the published EPSS model
documentation~\cite{firstepss} fixes the operational form agencies consume. The CISA KEV catalog
and the Binding Operational Directive program~\cite{cisakev,cisabod2201,cisabod2301} turned
exploited-in-the-wild evidence into binding remediation mandates and asset-visibility
requirements, which is what makes a finite remediation budget a hard operational constraint rather
than a soft preference. All of these signals are CVE-level and host-agnostic: they rank
vulnerabilities, not the hosts those vulnerabilities sit on. CAP-G is complementary to every one
of them, adding a host-mission dimension on top of the exploit signal rather than competing with
it, and inheriting their exploit estimates unchanged.

\textbf{Hygiene- and host-augmented prioritization.} The closest line of work augments exploit
scores with host-level hygiene, such as patch debt, directory exposure, and telemetry freshness,
raising precision on a hygiene-defined target by surfacing the badly-maintained hosts that an
exploit-only ranking overlooks~\cite{paper4hygieneprio}. The base scorer in this paper is exactly
that hygiene-augmented scorer (Section~\ref{sec:method}), reproduced unchanged so that CAP-G is
measured as a strict increment on an established baseline rather than against a degraded one. The
present work continues a line of context-aware prioritization for critical-infrastructure
and public-sector fleets~\cite{paper1vulnprio}: where the hygiene step moved the unit of analysis from the CVE to
the host's maintenance posture, CAP-G moves it again from maintenance posture to mission value.
Such hygiene scorers still weight the probability that a host is a weak link rather than the
consequence of its compromise, which is precisely the gap CAP-G addresses, and the homogeneous
control in Section~\ref{sec:results} demonstrates that the mission increment is separable from the
hygiene base rather than entangled with it.

\textbf{Risk- and asset-aware prioritization.} Federal guidance~\cite{nist80053,fips199,%
nist80060} and the CJIS policy~\cite{cjis2023} prescribe risk- and impact-driven remediation: the
Risk Management Framework requires that systems be categorized by mission impact and that
remediation follow that categorization rather than raw severity. But this guidance specifies an
objective, not a mechanism. It does not give a ranking algorithm, does not say how to combine an
impact category with an exploit probability, and does not quantify when asset weighting helps a
capacity-constrained queue and when it is inert. CAP-G fills that gap with a concrete, untuned
scorer and, more importantly, with a pre-registered and mechanism-controlled measurement of the
benefit, including the capacity regime in which it vanishes and the homogeneous case in which it
is provably zero. To our knowledge the capacity-dependence and the heterogeneity-dependence of the
asset-weighting benefit have not previously been isolated and bounded in this way,
and Section~\ref{sec:theory} supplements the measurement with a closed-form account
of why the multiplicative context factor reorders only near-tied pairs, why the
gain scales with the variance of the asset context score and vanishes exactly when
that variance does, and why the realized effect is bounded.

\textbf{Autonomic and self-managing security.} The broader framing for continuously reordering a
risk queue from changing host and exploit state is the autonomic-computing vision of self-managing
systems~\cite{kephart2003vision}, which anticipated security operations that monitor, prioritize,
and remediate with minimal human steering. CAP-G is one concrete scoring component such a system
would use to decide what to fix first, and our results say where in that loop asset context is
load-bearing and where it saturates.

\textbf{Honest evaluation.} Methodologically, the study follows the pre-registration discipline
now standard in empirical software engineering~\cite{wohlin2012experimentation}: hypotheses,
metrics, thresholds, and seeds are locked in a dated protocol before any evaluation-seed result is
inspected, failure criteria are explicit, and uncertainty is reported as bias-corrected and
accelerated bootstrap intervals~\cite{efron1987bca,efron1994introduction} with interval non-overlap
as the evidentiary standard rather than significance testing. The most visible consequence of that
posture is the willingness to report the primary hypothesis as partially supported rather than
re-cut the locked bar to claim a clean pass, which is intended to keep the evaluation credible
under adversarial review.
